\providecommand{\ZZ}{\mathbb{Z}}
\providecommand{\QQ}{\mathbb{Q}}
\providecommand{\CC}{\mathbb{C}}
\providecommand{\cO}{\mathcal{O}}
\providecommand{\Sp}{\mathrm{Sp}}
\providecommand{\wt}{\mathrm{wt}}
\providecommand{\Aut}{\mathrm{Aut}}
\providecommand{\CoHA}{\mathrm{CoHA}}
\providecommand{\BPS}{\mathrm{BPS}}
\providecommand{\CHL}{\mathrm{CHL}}
\providecommand{\DT}{\mathrm{DT}}
\providecommand{\kBKM}{\kappa_{\mathrm{BKM}}}
\providecommand{\kch}{\kappa_{\mathrm{ch}}}
\providecommand{\kcat}{\kappa_{\mathrm{cat}}}
\providecommand{\kfib}{\kappa_{\mathrm{fiber}}}
\providecommand{\DeltaFive}{\Delta_5}
\providecommand{\PhiTen}{\Phi_{10}^{\mathrm{un}}}
\providecommand{\PhiTenOP}{\Phi_{10}^{\mathrm{OP}}}
\providecommand{\Dfive}{D_5}
\providecommand{\ZOP}{Z_{\mathrm{OP}}}
\providecommand{\Yplus}{Y^+}
\providecommand{\Winfty}{\mathcal{W}_{1+\infty}}

\section*{Physical Meaning of CHL Denominators and Diagonal-Divisor Forms}
\label{wn:sec:physical-meaning-chl-dd-forms}

\paragraph{Theorem (Borcherds weight).}
Let \(R\) be a row label carrying a fixed Jacobi input \(\psi_R\), a
cusp normalisation, a multiplier, and a target paramodular group.  Write
\[
  c_R(0)=[q^0\zeta^0]\psi_R .
\]
Borcherds' singular theta lift gives the rowwise automorphic invariant
\[
  \kBKM(\Phi_R)=\wt(\Phi_R)=\frac{c_R(0)}{2}.
\]
The row label is part of the theorem.  Primitive \(\CHL\) denominators,
Nikulin order rows, and Gritsenko-Clery diagonal-divisor rows share the
row \(\DeltaFive\) and otherwise use different indexing principles.
A physical BPS interpretation is a second theorem, supplied after the
source geometry, scalar denominator, chamber, and Hall orientation have
been constructed.

\subsection*{Primitive CHL Denominator Ladder}

\paragraph{Theorem (primitive compact CHL rows).}
The primitive compact \(K3\times E\) \(\CHL\) denominator ladder is
indexed exactly by
\[
  N\in\{1,2,3,4,6\}.
\]
These are the symplectic K3 orders compatible with the elliptic factor
in the compact \(K3\times E\) source package.  Gritsenko-Nikulin, Gritsenko,
Eguchi-Hikami, and Govindarajan-Krishna give the Borcherds denominators
\(\Phi_N^{\CHL}\) with constants
\[
\begin{array}{c|ccccc}
N&1&2&3&4&6\\ \hline
c_N^{\CHL}(0)&10&8&6&4&2\\
\kBKM(\Phi_N^{\CHL})&5&4&3&2&1
\end{array}
\]
in the primitive denominator normalisation.  The \(N=1\) row is
\[
  \Phi_1^{\CHL}=\DeltaFive,\qquad
  \kBKM(\DeltaFive)=10/2=5.
\]
The \(N=5,7,8\) order rows belong to the Nikulin extension below.  A
common compact host and a row map are separate theorem data.

The \(N=1\) scalar inverse-square conventions separate as
\[
  \PhiTen=\DeltaFive^2,\qquad
  \Dfive=64^{-1}\DeltaFive,\qquad
  \PhiTenOP=\Dfive^2=4096^{-1}\DeltaFive^2,
\]
\[
  \ZOP^{K3\times E}=-(\PhiTenOP)^{-1}
  =-\Dfive^{-2}
  =-4096\,\DeltaFive^{-2}.
\]
The unnormalised Igusa square \(\PhiTen\) is the scalar Siegel form of
weight \(10\).  The OP reduced-DT scalar uses the monic product
\(\Dfive\).  The chiral half has Borcherds weight \(5\).

\subsection*{Nikulin Order Extension}

\paragraph{Automorphic theorem (Nikulin order rows).}
Nikulin's symplectic K3 orders include \(5,7,8\).  In the order-indexed
automorphic extension their constants are
\[
\begin{array}{c|ccc}
N&5&7&8\\ \hline
\text{cycle shape}&1^4 5^4&1^3 7^3&1^2 2^1 4^1 8^2\\
c_N^{\mathrm{Nik}}(0)&4&2&2\\
\kBKM(\Phi_N^{\mathrm{Nik}})&2&1&1
\end{array}
\]
where \(\Phi_N^{\mathrm{Nik}}\) denotes the order-indexed Nikulin
Borcherds product after its Jacobi input, Weyl chamber, Weyl vector,
and reflection character have been fixed.

The \(N=5\) geometric test object is a Borcea-Voisin threefold
\[
  \widehat X_5^{\mathrm{BV}}\longrightarrow
  (S_5\times E)/(\iota_S\times\iota_E),
\]
with \(S_5\) carrying an order-\(5\) symplectic automorphism commuting
with the anti-symplectic involution \(\iota_S\).  The automorphic value
is \(\kBKM(\Phi_5^{\mathrm{Nik}})=2\).  A compact BPS theorem at this
row requires a reduced obstruction theory, orientation data, the
weight-\(2\) Borcherds product, and Hall-Borcherds recognition.

At \(N=7\), the direct quotient
\((S_7\times E)/\langle g_7\times\mathrm{id}_E\rangle\) has three
elliptic \(A_6\)-surface singularity curves.  A free translation
quotient belongs to a different host problem.  The automorphic constant
remains \(c_7^{\mathrm{Nik}}(0)=2\), hence
\(\kBKM(\Phi_7^{\mathrm{Nik}})=1\).  Gerbe-valued or twisted
\(K\)-theoretic refinements require a constructed torsion differential
character and a compatible Hall orientation.

At \(N=8\), the order-indexed value
\(\kBKM(\Phi_8^{\mathrm{Nik}})=1\) is automorphic data.  A compact
Calabi-Yau threefold source package is additional structure.  Hyperkahler
or stacky boundary evidence gives automorphy at the boundary; a \(1/4\)-BPS
Hilbert-space index requires the compact source package and its
orientation data.

\subsection*{Gritsenko-Clery Diagonal-Divisor Catalogue}

\paragraph{Theorem (diagonal-divisor rows).}
Gritsenko-Clery classify the genus-two modular forms whose divisor is
the diagonal \(\mathcal H_1\) with vanishing order \(1\) on a
paramodular quotient.  The catalogue is indexed by triples \((t,N;k)\),
where \(k\) is the weight:
\[
 \mathsf{GC}_8=
 \{(1,1;5),(2,1;2),(1,2;3),(3,1;1),(1,3;2),
 (4,1;\tfrac12),(1,4;\tfrac32),(2,2;1)\}.
\]
In the canonical order the rows are
\[
\begin{array}{c|c|c|c|c}
R&(t,N;k)&F_R&c_R(0)&\kBKM(F_R)\\ \hline
1&(1,1;5)&\DeltaFive&10&5\\
2&(2,1;2)&\Delta_2&4&2\\
3&(1,2;3)&\nabla_3&6&3\\
4&(3,1;1)&\Delta_1&2&1\\
5&(1,3;2)&\nabla_2&4&2\\
6&(4,1;\tfrac12)&\Delta_{1/2}&1&\tfrac12\\
7&(1,4;\tfrac32)&\nabla_{3/2}&3&\tfrac32\\
8&(2,2;1)&Q_1&2&1
\end{array}
\]
so the weight sequence is
\[
  (5,2,3,1,2,\tfrac12,\tfrac32,1),
\]
and the twice-weight sequence is
\[
  (10,4,6,2,4,1,3,2).
\]
The half-integral entries carry multiplier systems of orders \(8\) and
\(4\), respectively.  The catalogue overlaps the primitive \(\CHL\)
ladder at the triple \((1,1;5)\), where \(F_R=\DeltaFive\).  Order alone
is insufficient for comparison with the primitive \(\CHL\) table.  Any
further comparison requires a row map specifying the paramodular level,
divisor, multiplier, Jacobi input, and scalar-square normalisation.

\subsection*{Dyon Counting}

\paragraph{Physical theorem (primitive CHL dyon count).}
For a compact primitive \(\CHL\) row with charges \((Q_e,Q_m)\), choose
chemical potentials \((\rho,\sigma,v)\) in the standard \(N=4\) quadratic
normalisation.  The protected dyon index is obtained from the physics
scalar denominator \(\Psi_N^{\mathrm{phys}}\) on the Sen contour:
\[
  d_N(Q_e,Q_m)=(-1)^{Q_e\cdot Q_m+1}
  \oint_{\mathcal C_N}
  \frac{\exp\!\bigl(-i\pi(\rho Q_e^2+\sigma Q_m^2+
  2v\,Q_e\cdot Q_m)\bigr)}
  {\Psi_N^{\mathrm{phys}}(\rho,\sigma,v)}
  \,d\rho\,d\sigma\,dv .
\]
At \(N=1\), \(\Psi_1^{\mathrm{phys}}=\PhiTen=\DeltaFive^2\), the
Dijkgraaf-Verlinde-Verlinde and Sen denominator for the \(1/4\)-BPS
dyon index.  For \(N>1\), the scalar denominator and chamber are those
of the corresponding Jatkar-Sen or Govindarajan-Krishna \(\CHL\)
model.  The primitive Borcherds weight \(c_N^{\CHL}(0)/2\) is the
chiral-half automorphic invariant attached to the same row.

\paragraph{Semiclassical consequence.}
With
\[
  \Delta(Q_e,Q_m)=Q_e^2Q_m^2-(Q_e\cdot Q_m)^2
\]
in the standard \(N=4\) charge convention, the attractor saddle for an
actual \(\CHL\) throat gives
\[
  S_{\mathrm{BH}}^{(N)}(Q_e,Q_m)
  =
  \pi\sqrt{\Delta(Q_e,Q_m)/N}
  +O(\log \Delta).
\]
The logarithmic coefficient is read from the full scalar denominator,
the massless spectrum, and the chosen chamber.  The Gritsenko-Clery
diagonal-divisor weights supply automorphic Borcherds weights for their
own rows.  Substitution into the \(\CHL\) black-hole entropy formula
requires a row-map theorem and a compact throat theorem.

\subsection*{Mathieu Twining}

\paragraph{Theorem (twined elliptic genera).}
Mathieu moonshine supplies twined K3 elliptic genera
\(\phi_{\mathrm{ell}}^{g}(\tau,z)\) for \(g\in M_{24}\).  A twined
genus becomes a Borcherds product after three choices are fixed:
the normalised Jacobi input, the cusp convention, and the target
paramodular group.  The cycle shape of \(g\) is a conjugacy-class
invariant; the Borcherds weight is computed from \(c_R(0)/2\).

For the primitive \(\CHL\) ladder the constants are
\[
  (c_N^{\CHL}(0))_{N=1,2,3,4,6}=(10,8,6,4,2).
\]
For the Nikulin order extension the constants are
\[
  (c_N^{\mathrm{Nik}}(0))_{N=5,7,8}=(4,2,2).
\]
For the Gritsenko-Clery diagonal-divisor catalogue the constants are
\[
  (10,4,6,2,4,1,3,2)
\]
in triple-indexed order.  These three lines use different row systems.

\subsection*{Hall Denominators and Chiral Scope}

\paragraph{Scope theorem.}
A compact Calabi-Yau threefold source enters the CY-to-chiral functor at
\(d=3\), hence its output \(A=\Phi_3(\mathcal C_X)\) is \(E_1\)-chiral.
An \(E_2\)-structure belongs to a derived centre, a representation
category, or another explicitly constructed higher-centre layer.

The Borcherds side is the Hall-Drinfeld denominator branch.  The
Mukai self-mirror K3-current branch is a different construction built
from K3 stable-envelope and Mukai-current data.  A comparison between a
compact Hall algebra and a Borcherds denominator algebra requires a
positive half, a Hall product, a non-degenerate pairing after completion,
PBW comparison, relation preservation, parity of root multiplicities,
and compatibility with the denominator identity.

The local toric reference point remains
\[
  \CoHA(\CC^3)=\Yplus .
\]
The algebra \(\Winfty\) appears at Drinfeld-double, centre, completion,
or evaluation layers.  This convention fixes the positive-half side of
any compact-source comparison.

\subsection*{Scalar Assignments on the Common \(K3\times E\) Row}

The common triple \((1,1;5)\) carries the \(K3\times E\) Borcherds
denominator \(\DeltaFive\).  The compact product and its
Heisenberg-Mukai specialisation carry the assignments
\[
  \kcat(K3\times E)=\chi(\cO_{K3})\chi(\cO_E)=2\cdot0=0,
\]
\[
  \kch(K3\times E)=
  \sum_{q=0}^{3}(-1)^q h^{0,q}(K3\times E)
  =1-1+1-1=0,
\]
\[
  \kch^{\mathrm{Heis}}(K3\times E)=3,\qquad
  \kBKM(\mathfrak g_{\DeltaFive})=5,\qquad
  \kfib=24.
\]
These assignments come from five constructions: the total-space
categorical Euler characteristic, the compact Hodge supertrace, the
specialised Heisenberg-Mukai chiral algebra, the Borcherds denominator
weight of \(\DeltaFive\), and the Mukai-lattice rank of the K3 fibre.
The automorphic rule is
\[
  \kBKM(\Phi_R)=c_R(0)/2
\]
for the chosen row \(R\).
The expression
\(\kBKM=\kch+\chi(\cO_{\mathrm{fiber}})\) gives \(0\) on the
\(\DeltaFive\) compact row while \(\kBKM(\DeltaFive)=5\); the mismatch
\(0\ne5\) rules it out as a row identity.

\subsection*{Primary Theorems Used}

Borcherds' singular theta lift gives
\(\mathrm{wt}\operatorname{Bor}(\varphi)=c(0)/2\).  Igusa and
Gritsenko-Nikulin identify the \(N=1\) Borcherds product with
\(\DeltaFive\), with \(c_1(0)=10\) and weight \(5\).  Gritsenko,
Eguchi-Hikami, Govindarajan-Krishna, Jatkar-Sen, and Sen fix the
primitive \(\CHL\) denominator rows and the scalar \(1/4\)-BPS
contours.  Gritsenko-Clery give the diagonal-divisor catalogue and its
twice-weight tuple \((10,4,6,2,4,1,3,2)\).  Oberdieck-Pandharipande
fix the OP reduced-DT scalar normalisation
\(\PhiTenOP=4096^{-1}\DeltaFive^2\) and
\(\ZOP=-4096\,\DeltaFive^{-2}\).  Kontsevich-Soibelman and
Davison-Meinhardt provide the Hall product and BPS Lie algebra
formalism after orientation and purity hypotheses are fixed.

\subsection*{Data Required for a Physical Row Theorem}

An automorphic row becomes a compact physical \(\BPS\) theorem after
the following data have been supplied:
\[
\begin{gathered}
  \text{a compact or stacky source geometry }X_R,\quad
  \text{a reduced obstruction theory and orientation},\\
  \text{a Jacobi input with specified cusp normalisation},\quad
  \text{a scalar denominator and contour chamber},\\
  \text{a Hall product, pairing, PBW comparison, and signed-root parity},\\
  \text{a proof that the protected index equals the chosen product.}
\end{gathered}
\]
For \(N=1\), Oberdieck-Pandharipande and the Igusa square give the
scalar \(K3\times E\) branch
\(-4096\,\DeltaFive^{-2}\).  Gritsenko-Nikulin and Borcherds give
the \(\DeltaFive\) BKM denominator.  For \(N=2,3,4,6\), the primitive
\(\CHL\) denominator constants are the table above.  For \(N=5,7,8\),
the Nikulin order constants are automorphic boundary data.  For the
Gritsenko-Clery catalogue, each triple-indexed diagonal-divisor row
needs its own source theorem before it has a compact BPS interpretation.
